\section{Sprint Reports}
\label{sec:rga}

The work is organised into three sprints that move from source discovery to a
verified SQLite database. Each report uses the same structure: overview,
backlog, goal, team contributions, completed activities, difficulties and
responses, and a checkable output. This makes the difference between planned
work and delivered evidence visible. Figure~\ref{fig:sprint-goal-matrix}
summarizes the ownership of G1--G11 before the detailed reports.

\begin{figure}[H]
\centering
\begin{tikzpicture}[
  matrix of nodes,
  nodes={draw=culine,minimum width=1.58cm,minimum height=0.62cm,
    align=center,font=\scriptsize},
  row 1/.style={nodes={fill=cublue!15,font=\scriptsize\bfseries}},
  column 1/.style={nodes={minimum width=2.55cm,font=\scriptsize\bfseries}}
]
\matrix (goals) {
Goal & Sprint 1 & Sprint 2 & Sprint 3 \\
G1 Collection & |[fill=cuteal!45]| Achieved & & \\
G2 Selection & |[fill=cuteal!45]| Achieved & & \\
G3 Extraction & & |[fill=cuteal!45]| Achieved & \\
G4 Cleaning & & |[fill=cuteal!45]| Achieved & \\
G5 ERD & & |[fill=cuamber!35]| Developed & |[fill=cuteal!45]| Finalized \\
G6 BCNF & & & |[fill=cuteal!45]| Achieved \\
G7 Schema & & & |[fill=cuteal!45]| Achieved \\
G8 ETL & & & |[fill=cuteal!45]| Achieved \\
G9 Integrity & & & |[fill=cuteal!45]| Achieved \\
G10 Queries & & & |[fill=cuteal!45]| Achieved \\
G11 Deployment & & & |[fill=cuteal!45]| Achieved \\
};
\end{tikzpicture}
\caption{Sprint-to-goal matrix for G1--G11. Blank cells indicate that a goal
was not assigned to that sprint.}
\label{fig:sprint-goal-matrix}
\end{figure}

\clearpage
\subsection{Sprint 1: Data Discovery and Acquisition}
\label{subsec:sprint1}

\subsubsection{Sprint Overview}

Sprint 1 establishes the source collection. Team members search official
Bangladesh organisations and recognized international data portals, then
compare the candidate material in a shared register. A publication is not
accepted merely because its title concerns the environment. The review checks
the publisher, file type, time range, measurement grain, identifiers and the
possibility of connecting the records to the approved scope.

The Drive review identifies \mSourceDataFiles{} source/data files and
\mSupportFiles{} supporting or administrative files. Only
\mSourceOrganisations{} organisations provide the \mSourceFiles{} structured
files selected for pipeline processing. The other
collected material remains part of the discovery record, with its exclusion
reason, rather than being described as loaded data.

\subsubsection{Sprint Backlog}

\begin{enumerate}[leftmargin=*]
  \item Search national and international organisations for relevant
  environmental datasets and publications.
  \item Store candidate files without overwriting the original downloads and
  record their publisher and source location.
  \item Inspect content, format, time coverage, observation grain and usable
  identifiers.
  \item Select material that fits the approved environment-sector scope and
  record why other files or sheets are excluded.
\end{enumerate}

\subsubsection{Sprint Goal}

The sprint addresses G1, Environmental Data Collection, and G2, Source
Verification and Selection. Completion requires both an organised collection
and a review trail showing why each retained source can move to extraction.
Figure~\ref{fig:sprint1-goals} connects the sprint directly to those goals.

\begin{figure}[H]
\centering
\begin{tikzpicture}[
  sprint/.style={draw=cublue,rounded corners=3pt,fill=cublue,text=white,
    minimum width=2.4cm,minimum height=1.05cm,align=center,font=\bfseries},
  goal/.style={draw=culine,rounded corners=3pt,fill=cugrey,text width=6.2cm,
    minimum height=0.9cm,align=center,font=\small},
  arr/.style={-{Stealth[length=2mm]},thick,draw=cuteal}
]
\node[sprint] (s1) {Sprint 1};
\node[goal,right=2.0cm of s1,yshift=0.65cm] (g1)
  {G1 - Environmental Data Collection};
\node[goal,right=2.0cm of s1,yshift=-0.65cm] (g2)
  {G2 - Source Verification and Selection};
\coordinate (split) at ($(s1.east)+(0.75cm,0)$);
\draw[arr] (s1.east)--(split)|-(g1.west);
\draw[arr] (split)|-(g2.west);
\end{tikzpicture}
\caption{Goals completed in Sprint 1.}
\label{fig:sprint1-goals}
\end{figure}

\clearpage
\subsubsection{Team Contributions}

Table~\ref{tab:sprint1-contributions} records principal ownership without
implying that a task was completed without group review.

\begin{table}[H]
\centering
\caption{Principal team contributions in Sprint 1}
\label{tab:sprint1-contributions}
\begin{tabular}{L{0.26\textwidth}L{0.63\textwidth}}
\toprule
Contribution area & Team members and contribution \\
\midrule
Dataset search and collection & Safwan Ahmed Rayed and Sheikh Tawsif Bin Aftab
searched for and collected relevant environmental publications. \\
Source documentation & Sheikh Tawsif Bin Aftab and Syed Mehraz Monwar recorded
source and dataset information in the shared register. \\
File organisation & Sheikh Tawsif Bin Aftab and Syed Mehraz Monwar organised
the collected files and supporting material in Google Drive. \\
File examination and selection & Sanjida Mahi reviewed the collected material
for content and suitability, with selection decisions discussed by the team. \\
Review and verification & Syed Mehraz Monwar checked the organisation and
documentation needed for later report evidence. \\
\bottomrule
\end{tabular}
\end{table}

\subsubsection{Activities Performed}

\begin{description}[style=nextline]
  \item[Independent source searching.] Members search separately across the
  assigned domain and bring the results into one register. This reduces the
  risk that the collection depends on one search path.
  \item[Official-source checking.] The review prioritizes responsible
  government agencies and recognized institutions. Secondary references can
  help locate a publication, but the original publisher and file are checked
  before selection.
  \item[Source recording.] The register preserves the organisation, material
  title, source location, file type, availability and selection status. The
  retained link lets a reviewer return to the publication.
  \item[Collection and organisation.] Downloaded files are grouped by source
  organisation in the shared Drive. Original material is preserved while
  extracted copies and later outputs are stored separately.
  \item[Scope and structure screening.] The team checks whether a file contains
  relevant measurements, whether its observation unit is clear and whether
  stable identifiers support integration.
  \item[Exclusion recording.] Narrative-only reports, repeated derivatives,
  unsupported national aggregates and files with unsuitable granularity or
  identifiers remain documented instead of being forced into the loader.
  \item[Handoff preparation.] Files selected for extraction are accompanied by
  their source metadata and initial notes about headers, dates, units and
  missing-value markers.
\end{description}

\subsubsection{Difficulties and Responses}

Table~\ref{tab:sprint1-challenges} records the principal difficulties and the
responses used during Sprint 1.

\begin{table}[H]
\centering
\caption{Sprint 1 difficulties and practical responses}
\label{tab:sprint1-challenges}
\small
\begin{tabular}{L{0.20\textwidth}L{0.28\textwidth}L{0.42\textwidth}}
\toprule
Difficulty & Effect on the work & Response \\
\midrule
Fragmented publication & Environmental material is distributed across many
organisations and portals. & Maintain one source register with organisation,
title, location, format and selection status. \\
Inaccessible or unstable links & Slow pages and broken downloads delay source
verification. & Search alternate pages belonging to the same publisher and
retain the working official location when found. \\
PDF-only information & Narrative pages or embedded tables cannot enter the
structured loader directly. & Prefer machine-readable files; retain relevant
reports as evidence or process them separately only when their grain is clear. \\
Inconsistent terminology & Similar measures and places appear with different
labels, units or spellings. & Record differences during discovery so they can
be reconciled explicitly in Sprint 2. \\
Scope mismatch & A file can be credible but too aggregated, duplicated or
outside the implemented domains. & Give it a documented exclusion reason while
keeping it in the source inventory. \\
\bottomrule
\end{tabular}
\end{table}

\subsubsection{Sprint Output}

Sprint 1 produces a source register and an organised collection containing
\mSourceDataFiles{} source/data files plus \mSupportFiles{} support files. The
screening outcome identifies \mSourceFiles{} structured files for the loader
and records excluded files and sheets separately. Section~\ref{sec:data}
provides the complete source-by-source accounting and shows that collected,
selected and loaded are different measurements.

\subsection{Sprint 2: Extraction and Data Preparation}
\label{subsec:sprint2}

\subsubsection{Sprint Overview}

Sprint 2 converts selected spreadsheets into explicit records. Each workbook
is inspected before automation because the files do not share one layout:
station names can appear above a block, dates can run across columns and one
sheet can contain several unrelated tables. The extraction scripts preserve
the original files and create raw 0NF CSV blocks for review.

The sprint also develops the integrated ERD. Modelling and extraction proceed
together because an entity, candidate key or relationship cannot be confirmed
until actual rows reveal the observation grain. The extraction ultimately
offers \mRowsOffered{} rows across \mSourceBlocks{} extracted source blocks;
the stage totals and changing record grain are explained in
Section~\ref{sec:etl}.

\subsubsection{Sprint Backlog}

\begin{enumerate}[leftmargin=*]
  \item Inspect the selected workbooks and locate their real data regions,
  repeated headers, dates, units and identifiers.
  \item Extract approved structures into raw, traceable 0NF blocks.
  \item Standardize missing markers, names, dates and safely parseable numeric
  formats while preserving uncertain cases for review.
  \item Record duplicates, conflicts, exclusions and source-precedence
  decisions.
  \item Develop source-level ERDs and integrate only relationships supported by
  dependable identifiers.
\end{enumerate}

\subsubsection{Sprint Goal}

Sprint 2 completes G3, Data Extraction, and G4, Data Cleaning and
Preprocessing. It develops G5, ERD Development and Integration, for final
review in Sprint 3.
Figure~\ref{fig:sprint2-goals} distinguishes completed goals from the goal that
continues into Sprint 3.

\begin{figure}[H]
\centering
\begin{tikzpicture}[
  sprint/.style={draw=cublue,rounded corners=3pt,fill=cublue,text=white,
    minimum width=2.4cm,minimum height=1.05cm,align=center,font=\bfseries},
  goal/.style={draw=culine,rounded corners=3pt,fill=cugrey,text width=6.2cm,
    minimum height=0.84cm,align=center,font=\small},
  arr/.style={-{Stealth[length=2mm]},thick,draw=cuteal}
]
\node[sprint] (s2) {Sprint 2};
\node[goal,right=2.0cm of s2,yshift=1.05cm] (g3) {G3 - Data Extraction};
\node[goal,right=2.0cm of s2] (g4) {G4 - Data Cleaning and Preprocessing};
\node[goal,right=2.0cm of s2,yshift=-1.05cm] (g5)
  {G5 - ERD Development and Integration};
\coordinate (split) at ($(s2.east)+(0.75cm,0)$);
\draw[arr] (s2.east)--(split)|-(g3.west);
\draw[arr] (split)--(g4.west);
\draw[arr] (split)|-(g5.west);
\end{tikzpicture}
\caption{Goals completed or developed in Sprint 2.}
\label{fig:sprint2-goals}
\end{figure}

\subsubsection{Team Contributions}

Table~\ref{tab:sprint2-contributions} records the principal contribution areas
for Sprint 2.

\begin{table}[H]
\centering
\caption{Principal team contributions in Sprint 2}
\label{tab:sprint2-contributions}
\begin{tabular}{L{0.26\textwidth}L{0.63\textwidth}}
\toprule
Contribution area & Team members and contribution \\
\midrule
File inspection & Sheikh Tawsif Bin Aftab and Syed Mehraz Monwar examined the
selected files to locate structures and information requiring extraction. \\
Extraction development & Sheikh Tawsif Bin Aftab and Syed Mehraz Monwar
developed and adjusted the Python extraction methods. \\
Data extraction and CSV preparation & Sheikh Tawsif Bin Aftab and Syed Mehraz
Monwar extracted the approved blocks and organised their raw CSV outputs. \\
Source-level ERD development & Safwan Ahmed Rayed and Sanjida Mahi developed
ERDs from the structures exposed by the selected sources. \\
ERD comparison and integration & Safwan Ahmed Rayed and Sanjida Mahi compared
entity grain, keys and relationship paths to develop the integrated model. \\
Review and verification & The group compared extracted samples with the source
files before the blocks moved to normalization. \\
\bottomrule
\end{tabular}
\end{table}

\subsubsection{Activities Performed}

\begin{description}[style=nextline]
  \item[Manual spreadsheet inspection.] The team identifies data regions,
  merged headings, daily and monthly wide layouts, repeated blocks and the
  cells that supply station, location, date and measurement context.
  \item[Source-structure identification.] Recurring patterns are described
  before coding. Particular care is given to labels placed outside the data
  rows because losing them would make the extracted observations anonymous.
  \item[Work distribution.] Selected workbooks are divided among members while
  common parsing and naming rules remain shared. Sample outputs are compared so
  parallel work does not create incompatible record layouts.
  \item[Python extraction.] Scripts based on \texttt{openpyxl} read approved
  sheets, carry contextual identifiers into each record and write explicit CSV
  blocks. Source-specific handling is retained where layouts genuinely differ.
  \item[Wide-to-long preparation.] Repeated month, day or measurement columns
  are exposed as values that can later become atomic 1NF rows. The original
  spreadsheet remains unchanged.
  \item[Cleaning preparation.] Utilities handle known missing markers, name
  mappings, numeric conversion and date interpretation. An uncertain value is
  logged rather than silently converted.
  \item[Exclusion and conflict tracking.] Files, sheets or blocks that cannot be
  integrated receive reasons. Duplicates, inconsistent totals and competing
  values remain visible in the quality records.
  \item[ERD integration.] Candidate entities are compared by meaning, key and
  grain. A link is included only when the sources provide a dependable joining
  identifier.
\end{description}

\subsubsection{Difficulties and Responses}

Table~\ref{tab:sprint2-challenges} connects the main Sprint 2 difficulties to
their practical responses.

\begin{table}[H]
\centering
\caption{Sprint 2 difficulties and practical responses}
\label{tab:sprint2-challenges}
\small
\begin{tabular}{L{0.20\textwidth}L{0.28\textwidth}L{0.42\textwidth}}
\toprule
Difficulty & Effect on the work & Response \\
\midrule
Context outside data rows & Station or district labels can appear only once
above a block. & Carry the contextual value into every extracted record in that
block. \\
Wide spreadsheet layouts & Months, days or measures spread across columns do
not form atomic database records. & Reshape them into row-based observations
while retaining source and block identity. \\
Merged and repeated headers & Generic readers can mistake headings for values
or omit part of the schema. & Inspect each workbook, define its true header
region and apply source-specific extraction rules. \\
Mixed date conventions & Year, fiscal year, month and day appear in different
positions and formats. & Parse them according to the block's declared grain and
keep unresolved dates out of the accepted output. \\
Inconsistent names and values & Spelling variants, missing markers and numeric
text obstruct joining and loading. & Apply documented mappings and parsing
rules; log uncertain conversions and source conflicts. \\
Large record volume & Manual copying is slow and difficult to reproduce. & Use
scripts to extract repeatably, then compare samples and counts with the source. \\
\bottomrule
\end{tabular}
\end{table}

\subsubsection{Sprint Output}

Sprint 2 produces the extraction scripts, quality and exclusion records,
\mSourceBlocks{} raw source blocks and \mRowsOffered{} rows offered to the
normalization workflow. It also produces the integrated ERD draft that Sprint
3 checks against the normalized relations. Sections~\ref{sec:etl} and
\ref{sec:dm} provide the detailed implementation and approved final ERD.

\subsection{Sprint 3: Normalization, Implementation and Verification}
\label{subsec:sprint3}

\subsubsection{Sprint Overview}

Sprint 3 turns the prepared blocks into the final normalized SQLite database.
Functional dependencies are checked from 1NF through BCNF, the integrated ERD
is finalized, the approved relational schema is implemented in SQLite and the data is
loaded through the reproducible build process. Verification then checks the
schema, relationships, queries, backup and restoration rather than treating a
successful load command as completion.

The verified database contains \mDbTables{} tables, \mDbViews{} views,
\mDbColumns{} columns, \mDbPk{} primary-key definitions,
\mDbFk{} foreign-key relationships and \mDbRows{} rows. Its declared time
coverage is \mYearMin--\mYearMax. These values are generated from the delivered
database rather than copied from a manually maintained summary.

\subsubsection{Sprint Backlog}

\begin{enumerate}[leftmargin=*]
  \item Normalize extracted relations through 1NF, 2NF, 3NF and BCNF while
  documenting functional dependencies and decompositions.
  \item Finalize the ERD and define primary, candidate and foreign keys in the
  final schema definition.
  \item Implement and load the schema in SQLite3 from the accepted normalized
  outputs.
  \item Validate required values, key uniqueness, constraints and foreign keys,
  and confirm that a clean rebuild contains the same database contents.
  \item Run competency questions and measure representative query execution
  times.
  \item Verify deployment, read-only operation, backup, restore and maintenance
  procedures.
\end{enumerate}

\subsubsection{Sprint Goal}

Sprint 3 finalizes G5 and completes G6--G11. The goal is not only a database
file but a project version whose schema, data and operating procedures can be
rebuilt and checked from the delivered project.
Figure~\ref{fig:sprint3-goals} shows the full set of goals completed in
Sprint 3.

\begin{figure}[H]
\centering
\begin{tikzpicture}[
  sprint/.style={draw=cublue,rounded corners=3pt,fill=cublue,text=white,
    minimum width=2.1cm,minimum height=1.0cm,align=center,font=\bfseries},
  goal/.style={draw=culine,rounded corners=3pt,fill=cugrey,text width=4.25cm,
    minimum height=0.72cm,align=center,font=\scriptsize},
  arr/.style={-{Stealth[length=1.8mm]},thick,draw=cuteal}
]
\node[sprint] (s3) {Sprint 3};
\node[goal,left=1.65cm of s3,yshift=1.25cm] (g5) {G5 - Final ERD};
\node[goal,left=1.65cm of s3,yshift=0.42cm] (g6) {G6 - Normalization to BCNF};
\node[goal,left=1.65cm of s3,yshift=-0.42cm] (g7) {G7 - Schema Implementation};
\node[goal,left=1.65cm of s3,yshift=-1.25cm] (g8) {G8 - Data Loading and ETL};
\node[goal,right=1.65cm of s3,yshift=1.25cm] (g9) {G9 - Integrity Checking};
\node[goal,right=1.65cm of s3,yshift=0.42cm] (g10) {G10 - Query and Demonstration};
\node[goal,right=1.65cm of s3,yshift=-0.42cm] (g11) {G11 - Deployment and Maintenance};
\coordinate (leftsplit) at ($(s3.west)+(-0.65cm,0)$);
\coordinate (rightsplit) at ($(s3.east)+(0.65cm,0)$);
\draw[arr] (s3.west)--(leftsplit)|-(g5.east);
\draw[arr] (leftsplit)|-(g6.east);
\draw[arr] (leftsplit)|-(g7.east);
\draw[arr] (leftsplit)|-(g8.east);
\draw[arr] (s3.east)--(rightsplit)|-(g9.west);
\draw[arr] (rightsplit)|-(g10.west);
\draw[arr] (rightsplit)|-(g11.west);
\end{tikzpicture}
\caption{Goals finalized or completed in Sprint 3.}
\label{fig:sprint3-goals}
\end{figure}

\subsubsection{Team Contributions}

Table~\ref{tab:sprint3-contributions} records the principal contribution areas
for Sprint 3.

\begin{table}[H]
\centering
\caption{Principal team contributions in Sprint 3}
\label{tab:sprint3-contributions}
\begin{tabular}{L{0.26\textwidth}L{0.63\textwidth}}
\toprule
Contribution area & Team members and contribution \\
\midrule
ERD refinement and normalization & Sanjida Mahi and Safwan Ahmed Rayed refined
the integrated ERD and worked through the normalization stages to BCNF. \\
Schema implementation and ETL & Sheikh Tawsif Bin Aftab implemented the final
schema, loading process and reproducible database build. \\
Database testing & Sheikh Tawsif Bin Aftab and Syed Mehraz Monwar checked the
implemented tables, relationships and query behaviour. \\
Review and verification & Syed Mehraz Monwar reviewed report evidence and
verified that explanations and diagrams agreed with the implemented database. \\
Queries and demonstrations & Sheikh Tawsif Bin Aftab and Syed Mehraz Monwar
prepared and reviewed saved read-only queries and database demonstrations. \\
Deployment and maintenance & Sheikh Tawsif Bin Aftab prepared the setup,
verification, backup and restoration workflow, with the results reviewed
by the team. \\
\bottomrule
\end{tabular}
\end{table}

\subsubsection{Activities Performed}

\begin{description}[style=nextline]
  \item[Relational-model refinement.] Extracted attributes are grouped by
  stable meaning and observation grain. Candidate keys and relationship paths
  are compared with the approved final ERD.
  \item[Normal-form evaluation.] Functional dependencies are checked at each
  stage. Partial and transitive dependencies are decomposed, then every BCNF
  determinant is checked against the candidate keys.
  \item[Schema implementation.] The final schema defines exact table
  names, columns, types, required values, checks, primary keys and foreign keys.
  The generated SQL and SQLite database are compared with that definition.
  \item[Normalization and ETL execution.] The pipeline produces intermediate
  files for review, loads the accepted BCNF files and records stage counts
  without assuming that changes in grain must be monotonic.
  \item[Database verification.] The database is checked for table and column
  agreement, primary-key uniqueness, required values, constraint behaviour and
  foreign-key consistency, and a clean rebuild is confirmed to contain the same
  database contents.
  \item[Query and demonstration preparation.] Saved read-only queries cover
  climate, wind, rivers, water quality, forestry, industry, relationship
  metadata, coverage, missing values and integrity.
  \item[Deployment and maintenance checks.] Linux and Windows procedures,
  read-only access, safe rebuild behaviour, backup and restore are exercised
  against temporary copies before final acceptance.
\end{description}

\subsubsection{Competency Query Validation}
\label{subsec:competency-query}

Competency questions test whether the schema can answer the environmental and
relationship questions implied by the final ERD. They therefore check more
than table existence. Table~\ref{tab:competency-questions} records the questions
and the verified Group-07 outcomes before representative SQL is shown.

\begin{table}[H]
\centering
\caption{Competency questions and verified Group-07 results}
\label{tab:competency-questions}
\begin{tabular}{C{0.07\textwidth}L{0.57\textwidth}L{0.25\textwidth}}
\toprule
ID & Question & Verified result \\
\midrule
CQ1 & How many temperature records are available for 2020? &
\mTemperatureRowsTwentyTwenty{} rows \\
CQ2 & How many districts have retained forest-area records? & 35 districts \\
CQ3 & Can every loaded water-quality row reach its retained river or water-body
reference through its station? &
\mWaterQualityConnected{} joined rows \\
CQ4 & Which adequately sampled station has the highest average monthly
rainfall? & \mWettestStation{}, \mWettestAverage{} mm \\
CQ5 & Are the retained industry-use observations retrievable by fiscal period?
& \mIndustryRows{} rows \\
CQ6 & What is the database's declared time coverage? &
\mYearMin--\mYearMax \\
CQ7 & Which geographic limitation remains visible? &
\mUnmappedStations{} stations lack a supported district mapping \\
\bottomrule
\end{tabular}
\end{table}

The following query checks the Station--Temperature relationship and uses the
implemented measurement column, \texttt{Temp}.

\begin{lstlisting}[language=SQL,caption={Station and temperature relationship validation},
label={lst:station-temperature-validation}]
SELECT s.Station_Name, t.Year, t.Month, t.Type, t.Temp
FROM Station AS s
JOIN Temperature_Record AS t
  ON s.Station_Name = t.Station_Name
WHERE t.Year = 2020;
\end{lstlisting}

The next query demonstrates the relationship tested by CQ3. The joining column
is \texttt{WQ\_Station\_Name} in both implemented relations.

\begin{lstlisting}[language=SQL,caption={Water-quality records joined to their river or water-body references},
label={lst:water-quality-validation}]
SELECT rs.River_Name,
       w.WQ_Station_Name,
       w.Year,
       w.Parameter_Type,
       w.Value
FROM Water_Quality AS w
JOIN River_Station AS rs
  ON rs.WQ_Station_Name = w.WQ_Station_Name;
\end{lstlisting}

CQ4 applies a minimum sample-size rule so a station represented by only a few
months cannot outrank a long series without disclosure.

\begin{lstlisting}[language=SQL,caption={Rainfall competency ranking},
label={lst:rainfall-ranking}]
SELECT Station_Name, COUNT(*) AS Measurements,
       ROUND(AVG(Rainfall), 2) AS Average_Rainfall,
       ROUND(MAX(Rainfall), 2) AS Highest_Rainfall
FROM Rainfall_Record
WHERE Rainfall IS NOT NULL
GROUP BY Station_Name
HAVING COUNT(*) >= 100
ORDER BY Average_Rainfall DESC
LIMIT 15;
\end{lstlisting}

Every saved query is warmed once and then executed \mBenchmarkRuns{} times
against SQLite \mSQLiteVersion{} in read-only mode. Table~\ref{tab:competency-timing}
reports median and 95th-percentile wall-clock time. These figures describe the
test machine and are not general performance guarantees.

\begin{table}[H]
\centering
\caption{Selected competency-query timing results}
\label{tab:competency-timing}
\begin{tabular}{L{0.46\textwidth}r r r}
\toprule
Check & Returned rows & Median (ms) & P95 (ms) \\
\midrule
\inputrows{scripts/generated/competency_results.tex}
\bottomrule
\end{tabular}
\end{table}

\begin{figure}[H]
\centering
\begin{tikzpicture}[
  box/.style={draw=cublue,rounded corners,fill=cugrey,text width=3.0cm,
    minimum height=0.95cm,align=center,font=\small},
  decision/.style={diamond,draw=cuteal,aspect=2.1,align=center,font=\small},
  arr/.style={-{Stealth[length=2mm]},thick,draw=cuteal},
  node distance=0.65cm
]
\node[box] (q) {Run competency queries};
\node[box,right=of q] (i) {Run integrity and FK checks};
\node[decision,right=of i] (pass) {All pass?};
\node[box,below=0.8cm of pass] (release) {Checks pass};
\node[box,below=0.8cm of i] (fix) {Correct data, schema or query};
\draw[arr] (q)--(i);
\draw[arr] (i)--(pass);
\draw[arr] (pass)--node[right,font=\scriptsize]{yes}(release);
\draw[arr] (pass)--node[above,font=\scriptsize]{no}(fix);
\draw[arr] (fix)-|(q);
\end{tikzpicture}
\caption{Competency-query and database-verification workflow.}
\label{fig:competency-query-flow}
\end{figure}

Figure~\ref{fig:competency-query-flow} shows that a failed competency,
integrity or foreign-key result returns the affected data, schema or query to
correction. Acceptance follows only after the questions and engine checks pass.

The peer group could not provide an executable SQLite database. Consequently,
the competency evaluation covers Group-07's verified implementation only and
does not present a two-database comparison. The separately received MySQL
demonstration DDL is not treated as an implemented peer database.

\subsubsection{Difficulties and Responses}

Table~\ref{tab:sprint3-challenges} records the implementation and verification
difficulties addressed during Sprint 3.

\begin{table}[H]
\centering
\caption{Sprint 3 difficulties and practical responses}
\label{tab:sprint3-challenges}
\small
\begin{tabular}{L{0.20\textwidth}L{0.28\textwidth}L{0.42\textwidth}}
\toprule
Difficulty & Effect on the work & Response \\
\midrule
Unmatched water-body names & Some water-quality stations do not match the river
register by a dependable name. & Retain verified station--river mappings in
\texttt{River\_Station}; do not invent geographic relationships for unmatched
station records. \\
Invalid calendar combinations & Separate year, month and day values can form
dates such as 30 February. & Use the time relations and checks so invalid
calendar combinations are rejected during mutation tests. \\
Malformed numeric text & Typographical decimal points and formatted numbers can
break numeric loading or change meaning. & Flag the source problem, preserve
its evidence and convert only when the intended value is defensible. \\
Conflicting totals & A published total can disagree with its components. & Keep
the published value distinguishable, calculate comparison totals separately
and record the discrepancy instead of silently replacing it. \\
Schema mismatch & A diagram, schema definition, generated SQL and database can disagree
after revision. & Regenerate and compare all four representations, then run the
schema, integrity and competency tests before acceptance. \\
Rebuild and restore risk & A script could overwrite the wrong database or
produce a structurally different copy. & Refuse unrelated destinations and
compare content checksums after clean rebuild, backup and restore. \\
\bottomrule
\end{tabular}
\end{table}

\subsubsection{Sprint Output}

Sprint 3 produces the approved final ERD, \mBCNFRels{} BCNF relations, the
\mDbRows{}-row SQLite database, \mDbViews{} read-only summary views, the schema
definition, saved queries, measured competency results, automated verification
tests and backup--restore utilities. The database passes
\texttt{PRAGMA integrity\_check}, returns \mDbFkViolations{} foreign-key
violations and rebuilds to the same content checksum. Section~\ref{sec:deploy}
summarizes these final deployment checks without duplicating
the full query evidence.
